\documentclass[12pt]{article}
\usepackage[utf8]{inputenc}
\usepackage{amsmath}
\usepackage{graphicx}
\usepackage{natbib}
\usepackage{hyperref}
\usepackage{lipsum} % For dummy text, you can remove this in your final document

\title{Synthetic Biology and Environmental Sustainability: Bioengineering Approaches for Waste Management, Carbon Capture, and Biodiversity Restoration}
\author{Your Name}
\date{\today}

\begin{document}

\maketitle

\begin{abstract}
Synthetic biology offers innovative approaches to address environmental challenges such as waste management, carbon capture, and biodiversity restoration. This paper reviews current applications of synthetic biology in these areas, discusses their benefits, environmental impacts, and examines future potential for advancing sustainability goals.
\end{abstract}

\section{Introduction}
Synthetic biology integrates engineering principles into biological systems, enabling the design and construction of novel biological entities for various applications. This paper explores the role of synthetic biology in promoting environmental sustainability, with a focus on bioengineering solutions for waste management, carbon capture, and biodiversity restoration.

\section{Applications Overview}
\subsection{Waste Management}
Utilization of synthetic biology for bioremediation of pollutants, plastic degradation, and wastewater treatment \citep{dvoøák2020microbial, turner2020microbial}.

\subsection{Carbon Capture}
Bioengineering approaches for enhancing carbon sequestration by engineered microorganisms and plants \citep{carrasco2020engineering, mcginn2020photosynthetic}.

\subsection{Biodiversity Restoration}
Synthetic biology applications in restoring degraded ecosystems, species conservation, and genetic rescue efforts \citep{debus2017synthetic, seddon2021genomic}.

\section{Case Studies}
\subsection{Bioengineering for Plastic Degradation}
Case studies on engineered microbes capable of breaking down plastics and their environmental implications \citep{wei2020engineering, gajendiran2021plastic}.

\subsection{Carbon-Neutral Biofuels}
Examples of synthetic biology advancements in producing sustainable biofuels from renewable resources \citep{paterlini2021synthetic, verma2016synthetic}.

\subsection{Genetic Rescue and De-extinction}
Ethical and ecological considerations in using synthetic biology for genetic rescue and de-extinction projects \citep{prowse2017should, vince2018wild}.

\section{Benefits}
\subsection{Environmental Impact Reduction}
Reducing environmental pollution and carbon footprint through sustainable biotechnological solutions \citep{chuang2020review, enrique2021role}.

\subsection{Innovative Solutions}
Innovative approaches to complex environmental challenges not easily addressed by traditional methods \citep{purnick2009design, ratner2020synthetic}.

\section{Challenges}
\subsection{Ethical Concerns}
Ethical implications of manipulating natural systems and biodiversity conservation \citep{redford2015synthetic, keulartz2013ecology}.

\subsection{Regulatory Frameworks}
Regulatory challenges and considerations in deploying synthetic biology solutions in the environment \citep{garfinkel2017synthetic, magner2020synthetic}.

\section*{References}
\bibliographystyle{plainnat}
\bibliography{prompt16_35}

\end{document}
